\section{Equivalences and an exhaustive seed cover}
\label{sec:seed}

The goal is to replace every possible second basis by one of finitely
many nine-phase seed balls. Closed phase discretization and necessary
compatibility tests follow the approach of
\cite[Section~3]{JamingEtAl2009} and
\cite[Sections~2--3]{JamingMatolcsiMora2010}. The symmetry reduction and
shared-coordinate estimates below prove completeness of this particular
seed cover.

We use phases in cycles and write
\[
 e(t)=\exp(2\pi i t),\qquad t\in\mathbb R,
\]
with entrywise application to vectors and matrices.  Matrix indices run from
$0$ to $5$, and a phase matrix is stored in row-major order, with $(i,j)$
in position $6i+j$.  The Hermitian inner product is
$\langle u,v\rangle=\sum_j\overline{u_j}v_j$, and $A^*$ denotes conjugate
transpose.  Vectors have their Euclidean norms; a matrix regarded as a phase
or displacement vector has Frobenius norm $\|\cdot\|_F$.  The notation
$\|\cdot\|_{\mathrm{op}}$ denotes the induced Euclidean operator norm.
We identify a complex equation with its real and imaginary parts, in that
order, so that a complex scalar and its realification have the same norm.
An angle is in radians only when this is stated explicitly.

\subsection{Equivalences and phase gauges}

\begin{definition}
\label{seed:extendibility}
A matrix $H\in\mathbb C^{6\times6}$ is a complex Hadamard matrix here if
every entry has modulus one and $HH^*=6I$.  It is \emph{extendible} if the
standard basis and the columns of $H/\sqrt6$ belong to a family of four
mutually unbiased orthonormal bases.
\end{definition}

A common unitary maps the first member of any proposed family to the
standard basis.  Unbiasedness to that basis forces the entries of the second
basis matrix to have modulus $1/\sqrt6$.  Thus it is $H/\sqrt6$ for a matrix
as in Definition~\ref{seed:extendibility}; since $H$ is square, $HH^*=6I$
also gives $H^*H=6I$.  Every such $H$ admits a real phase lift $P$ with
$H=e(P)$.

\begin{lemma}[Equivalences]
\label{seed:equivalences}
Extendibility is invariant under row or column permutations, multiplication
of rows or columns by unit scalars, complex conjugation, and transposition.
\end{lemma}

\begin{proof}
Permuting or rephasing columns only changes the order or phases of vectors
within the second basis.  A row permutation or row rephasing is a common
unitary operation on all four bases; it preserves the standard basis up
to the same irrelevant changes.  Complex conjugation preserves both
orthogonality and the moduli of inner products.

For transposition, put $U=H/\sqrt6$.  Apply the common unitary $U^*$ to a
family whose first two basis matrices are $I,U$.  These two matrices become
$U^*,I$.  Interchange their roles and conjugate every coordinate.  The first
two matrices are now $I,U^T$.  All these operations are reversible, which
proves the assertion in both directions.
\end{proof}

The transposition argument applies to an exact Hadamard matrix.  It does
not assert that an approximate center supplies a unitary change of basis.

For a real $6$-by-$6$ matrix $z$, define double centering by
\begin{equation}
\label{seed:projection-formula}
 (\Pi z)_{ij}
 =z_{ij}-\frac16\sum_bz_{ib}-\frac16\sum_az_{aj}
       +\frac1{36}\sum_{a,b}z_{ab}.
\end{equation}
Its range is the \emph{horizontal} subspace of matrices with every row sum
and every column sum zero.  We also use unwrapped dephasing,
\begin{equation}
\label{seed:dephasing-formula}
 D_0(P)_{ij}=P_{ij}-P_{i0}-P_{0j}+P_{00}.
\end{equation}

\begin{lemma}[Horizontal phase representatives]
\label{seed:horizontal}
The map $\Pi$ is the orthogonal projection onto the horizontal subspace.
In particular, $\|\Pi z\|_F\le\|z\|_F$.  If two phase representatives have
difference $z$, the second is equivalent under row and column rephasings
to the first plus $\Pi z$.  This statement remains valid after arbitrary
entrywise integer changes of phase representatives.
\end{lemma}

\begin{proof}
Formula~\eqref{seed:projection-formula} has zero row and column sums, and
$z-\Pi z$ is a sum of a matrix constant on each row and a matrix constant
on each column.  Every such sum is orthogonal to the horizontal subspace
in the real Frobenius inner product.  This proves the projection assertion
and its norm bound.  If $Q=P+z$, then
$P+\Pi z=Q-(z-\Pi z)$ differs from $Q$ by row and column phases, which are
allowed by Lemma~\ref{seed:equivalences}.  Entrywise integer additions
leave $e(P)$ unchanged.
\end{proof}

Likewise, $D_0(P)$ gives an equivalent matrix with first row and column
phases zero.  When a wrapped difference is used, its chosen integers define
legitimate phase representatives.  No argument requires a rounded decision
at a half-integer to agree with a uniquely preferred nearest representative.

\subsection{A symmetry arrangement for nine phases}

The \emph{gap} of a finite circular multiset is its largest circular gap.
Equivalently, it is the greatest length of an open arc containing no point
of the multiset.  Points are counted with multiplicity when sorted; a
repeated point contributes a zero gap.  For each pair of distinct rows
of $P$, consider their six phase differences modulo one, and do the same
for each pair of distinct columns.

\begin{lemma}[Seed arrangement]
\label{seed:arrangement}
Every exact Hadamard matrix has an equivalent dephased representative for
which a pair of rows attaining the minimum gap among all row and column
pairs is the first pair, and its phases can be written
\begin{equation}
\label{seed:x-arrangement}
 x=(0,x_1,x_2,x_3,x_4,x_5),\qquad
 0\le x_1\le\cdots\le x_5\le1,\qquad
 g_x=1-x_5,\qquad x_1\le x_5-x_4.
\end{equation}
After any one of these five noninitial columns is selected and moved to
second position, the last four rows can be permuted so that the second
column has phases
\begin{equation}
\label{seed:y-arrangement}
 y=(0,\xi,y_2,y_3,y_4,y_5),\qquad
 0\le y_2\le y_3\le y_4\le y_5\le1,
\end{equation}
where $\xi$ is the selected entry of $x$.  These lists satisfy
\begin{equation}
\label{seed:local-conditions}
 \sum_{j=0}^5 e(x_j)=0,\qquad
 \sum_{j=0}^5 e(y_j)=0,\qquad g_y\ge g_x.
\end{equation}
\end{lemma}

\begin{proof}
There are finitely many row and column pairs, so a pair of minimum gap
exists.  Transpose if it is a column pair, and permute rows to make it the
first pair.  Choose one of its largest gaps, start the circular ordering
immediately after that gap, and use this starting point as the first
column.  Dephase the first row and column.  Ordering the remaining columns
along the complementary arc gives $g_x=1-x_5$ and the sorted list in
\eqref{seed:x-arrangement}.

If the final orientation inequality fails, interchange the two selected
rows, use the old $x_5$ column as the starting column, and reverse the
ordering.  The new list is
\[
 x'_0=0,\qquad x'_i=x_5-x_{5-i}\quad(1\le i\le5).
\]
It has the same distinguished largest gap, and
$x'_1=x_5-x_4$ whereas $x'_5-x'_4=x_1$.  Thus this orientation satisfies
the required inequality whenever the former one does not.  The construction
uses a chosen largest gap and is unchanged in validity by ties or repeated
points.

Select a noninitial column and place it second.  Permuting the last four
rows sorts its last four phase entries, without changing either of the first
two rows.  Orthogonality of those rows gives the first vanishing sum in
\eqref{seed:local-conditions}; orthogonality of the first two columns gives
the second.  Row and column rephasings rotate the relevant difference
multisets; permutations merely reorder points or relabel pairs.  Transposition
interchanges the two collections of pairs.  Consequently their minimum
gap remains $g_x$, and the newly selected column pair has $g_y\ge g_x$.
\end{proof}

Set $N=112$ and $w=1/N$.  All bins are closed:
\begin{equation}
\label{seed:bins}
 I_b=[b/N,(b+1)/N],\qquad
 u_b=(b+1/2)/N,\qquad 0\le b<N.
\end{equation}
They cover $[0,1]$, including both endpoints.  Assignments at shared
boundaries need not be unique.  Sorted exact phases admit nondecreasing
bin indices, so write the five indices for $x_1,\ldots,x_5$ as
\[
 b^X=(a,b,c,d,l),\qquad 0\le a\le b\le c\le d\le l<N.
\]
The orientation condition implies the necessary integral restriction
\begin{equation}
\label{seed:orientation-filter}
 \frac aN\le x_1\le x_5-x_4\le\frac{l+1-d}{N},
 \qquad a\le l-d+1.
\end{equation}

For this index tuple choose $s\in\{0,\ldots,4\}$ as the first minimizer of
$|2b^X_s+1-N|$.  This is the rule in \texttt{seeds::mk}.  Its selected
original column is $s+1$, its selected exact phase is $\xi=x_{s+1}$, and
the new column order is the original column zero, column $s+1$, and the
other four columns in their former order.  This selection only specifies
a column; it imposes no further condition on $\xi$.

After arranging $y$ as in \eqref{seed:y-arrangement}, its five varying
entries have bins
\[
 b^Y=(b^X_s,b^Y_1,b^Y_2,b^Y_3,b^Y_4),\qquad
 0\le b^Y_1\le b^Y_2\le b^Y_3\le b^Y_4<N.
\]
The shared entry uses precisely the same bin as in $x$; it is not required
to be ordered relative to the four tail bins.

\begin{lemma}[Stability of circular gaps]
\label{seed:gap-stability}
If each point of one circular multiset is matched to a point of another
at circular distance at most $\delta$, their gaps differ by at most
$2\delta$.
\end{lemma}

\begin{proof}
Take an empty open arc of length $L$ in the first multiset.  If $L>2\delta$,
trim $\delta$ from each end.  A point of the second multiset inside the
remaining open arc would have its matched first point in the original
empty arc, a contradiction.  Hence the second gap is at least
$L-2\delta$.  If $L\le2\delta$, the same lower bound follows from
nonnegativity.  Apply this to a largest gap and then interchange the
two multisets.  No preservation of the sorted order is used.
\end{proof}

\begin{lemma}[Necessary midpoint gap filters]
\label{seed:gap-filters}
For every bin pair containing lists as in Lemma~\ref{seed:arrangement},
the midpoint gaps, including the fixed zero, satisfy
\begin{equation}
\label{seed:gap-inequalities}
 g_{x,\mathrm{mid}}\le1-\frac{l-1}{N},\qquad
 g_{y,\mathrm{mid}}\ge1-\frac{l+2}{N}.
\end{equation}
\end{lemma}

\begin{proof}
Each varying phase is within $w/2$ of its midpoint, and the initial zero
is unchanged.  Lemma~\ref{seed:gap-stability} bounds each gap change by $w$.
Since $x_5\in[l/N,(l+1)/N]$, we have
$1-(l+1)/N\le g_x\le1-l/N$.  Combining these inequalities with
$g_y\ge g_x$ proves both assertions.
\end{proof}

For sorted midpoint bin indices $q_0,\ldots,q_4$, the gap used by
\texttt{mk} is
\[
 \max\left\{\frac{q_0+1/2}{N},
 \frac{q_1-q_0}{N},\ldots,\frac{q_4-q_3}{N},
 1-\frac{q_4+1/2}{N}\right\}.
\]
For a $Y$ tuple it sorts a copy solely for this computation, preserving
the stored shared-first ordering.

\subsection{Vanishing-sum filters and the shared displacement}

For either list, let $u=(u_0,\ldots,u_4)$ be its five varying midpoints.
Define the real two-vector $c$ and real $2$-by-$5$ matrix $J$ by
\begin{equation}
\label{seed:residual-data}
 c=\frac{1+\sum_{j=0}^4e(u_j)}{2\pi i},\qquad J_j=e(u_j),
\end{equation}
using realification in both definitions.

\begin{lemma}[Necessary sum test]
\label{seed:sum-test}
If $1+\sum_{j=0}^4e(u_j+\Delta_j)=0$ and $|\Delta_j|\le w/2$, then
\begin{equation}
\label{seed:taylor-residual}
 \|c+J\Delta\|\le\frac{5\pi w^2}{4}.
\end{equation}
In particular, for every $v\in\mathbb R^2$,
\begin{equation}
\label{seed:sum-support}
 |v\cdot c|\le\frac{5\pi w^2}{4}\|v\|
                  +\frac w2\sum_{j=0}^4|v\cdot J_j|.
\end{equation}
\end{lemma}

\begin{proof}
The second derivative of $e(t)$ has modulus $4\pi^2$.  Taylor's formula
with integral remainder therefore gives
\[
 \left|\frac{e(u_j+\Delta_j)-e(u_j)}{2\pi i}
             -e(u_j)\Delta_j\right|\le\pi\Delta_j^2.
\]
Sum the five remainders and use the vanishing-sum hypothesis to obtain
\eqref{seed:taylor-residual}.  Taking the dot product with $v$, followed
by the triangle inequality and $|\Delta_j|\le w/2$, gives
\eqref{seed:sum-support}.
\end{proof}

The function \texttt{sumok} tests \eqref{seed:sum-support} in the residual
direction and in directions perpendicular to the five unit-exponential
columns.  The inequality holds for every real direction, so their numerical
construction need not reproduce any prescribed exact direction.  The
one-sided margins for evaluating the residuals and supports are established
in Section~\ref{sec:rounding}.

In grid units $h=N\Delta$, write
\begin{equation}
\label{seed:grid-error}
 E=\frac{5\pi}{4N},\qquad \|Jh+Nc\|\le E,
 \qquad |h_j|\le\frac12.
\end{equation}

\begin{lemma}[Shared-coordinate intervals]
\label{seed:shared-interval}
Let $t$ be the position of a shared entry in one of the two lists.  Every
feasible displacement satisfies, for each $v\in\mathbb R^2$,
\begin{equation}
\label{seed:clip-inequality}
 |Nv\cdot c+(v\cdot J_t)h_t|
 \le E\|v\|+\frac12\sum_{j\ne t}|v\cdot J_j|.
\end{equation}
If $a=v\cdot J_t\ne0$ and the right side is $r$, this is the closed
interval constraint
\begin{equation}
\label{seed:clip-interval}
 -\frac{Nv\cdot c}{a}-\frac r{|a|}
 \le h_t\le
 -\frac{Nv\cdot c}{a}+\frac r{|a|}.
\end{equation}
The intervals obtained for the shared entry from $X$ and $Y$ must intersect.
\end{lemma}

\begin{proof}
Separate the shared term in \eqref{seed:grid-error}, take its dot product
with $v$, and bound each remaining coordinate by $1/2$.  Solving the
resulting absolute-value inequality gives \eqref{seed:clip-interval} for
either sign of $a$.  The shared midpoint is identical in both lists, so
the same exact grid displacement belongs to both sets of necessary
intervals.
\end{proof}

The function \texttt{extend} begins with $[-1/2,1/2]$ and intersects these
constraints, using the shared column direction and directions perpendicular
to the other four columns.  Here $t=s$ for $X$ and $t=0$ for $Y$.
It omits divisors of magnitude below $10^{-6}$, which only removes a
constraint.  Numerator and endpoint margins enlarge the computed intervals,
as proved in Section~\ref{sec:rounding}.  The pair is rejected only when
the two resulting intervals are strictly disjoint; a common endpoint is
retained.

The enumeration in \texttt{seeds::build} now has an explicit interpretation.
It visits every weakly increasing five-bin $X$ tuple and applies
\eqref{seed:orientation-filter}, the first inequality in
\eqref{seed:gap-inequalities}, and \texttt{sumok}.  Separately, for every
possible first bin it visits all weakly increasing four-bin $Y$ tails and
applies \texttt{sumok}.  The program \texttt{covout} matches the shared
bin, applies the second gap inequality, and intersects the shared
displacement intervals.  Every discarded exact solution would therefore
violate one of the necessary conditions just proved.  The numerical gap
comparisons include outward slack, treated in Section~\ref{sec:rounding}.

\subsection{Certified reuse of seed centers}

The nine seed coordinates, before moving the selected column, are
\begin{equation}
\label{seed:nine-coordinates}
 q=(x_1,x_2,x_3,x_4,x_5,y_2,y_3,y_4,y_5)\in\mathbb R^9.
\end{equation}
For a bin pair let $u\in\mathbb R^9$ be its rational midpoint and put
$h=N(q-u)$.  Thus $h\in[-1/2,1/2]^9$; the shared displacement is $h_s$.
Retaining $u$ as a center would cover the entire bin with radius
$3/(2N)<0.032$.  To reduce the number of centers, we instead test whether
its feasible portion lies in a radius-$0.032$ ball about a previously
retained midpoint.  If no such test succeeds, we retain $u$.

Let $J_X,J_Y,c_X,c_Y$ denote the data in
\eqref{seed:residual-data}.  Define a real $4$-by-$9$ matrix $J_*$ by
specifying its action:
\begin{equation}
\label{seed:stacked-system}
 J_*h=\begin{pmatrix}
 J_X(h_0,h_1,h_2,h_3,h_4)^T\\
 J_Y(h_s,h_5,h_6,h_7,h_8)^T
 \end{pmatrix},\qquad
 p_*=\begin{pmatrix}-Nc_X\\-Nc_Y\end{pmatrix}.
\end{equation}
For feasible $h$ we have $J_*h=p_*+r_*$, where the two real two-vector
blocks $r_X,r_Y$ of $r_*$ satisfy $\|r_X\|,\|r_Y\|\le E$.

\begin{lemma}[Support bounds for center reuse]
\label{seed:dual-support}
Let $u'$ be a previously retained midpoint with the same selected position
$s$, and let $k=N(u-u')\in\mathbb Z^9$.  Every feasible displacement in
the current bin pair satisfies
\begin{equation}
\label{seed:distance-bound}
 N^2\|q-u'\|^2=\|k+h\|^2
 \le\|k\|^2+\frac94+2k\cdot h.
\end{equation}
Two upper bounds for the last term are
\begin{align}
 2k\cdot h&\le\sum_{j=0}^8|k_j|,
 \label{seed:box-support}\\
 2k\cdot h&\le
 \sum_{j=0}^8|k_j-(J_*^Tv)_j|
 +2v\cdot p_*+2E(\|v_X\|+\|v_Y\|),
 \qquad v\in\mathbb R^4,
 \label{seed:dual-formula}
\end{align}
where $v_X,v_Y$ are the two real two-vector blocks of $v$.
\end{lemma}

\begin{proof}
Expand the square in \eqref{seed:distance-bound} and use
$\|h\|^2\le9/4$.  The coordinate bounds give
\eqref{seed:box-support}.  For the other bound, write
\[
 2k\cdot h
 =2(k-J_*^Tv)\cdot h+2v\cdot p_*+2v\cdot r_*.
\]
The first term is at most $\sum_j|k_j-(J_*^Tv)_j|$ by the same coordinate
bounds.  Cauchy--Schwarz on the two residual blocks bounds the last term
by $2E(\|v_X\|+\|v_Y\|)$.
\end{proof}

The function \texttt{supdual} proposes vectors $v$ and evaluates the right
side of \eqref{seed:dual-formula}.  The validity of a proposal does not
depend on how it was generated or on the accuracy of the heuristic linear
solve.  Only finite proposed components of magnitude at most $10^4$ are
used.  Its evaluation replaces $1.25\pi/N$ by the larger $1.251\pi/N$
and adds a final support margin.  Section~\ref{sec:rounding} proves that
the resulting bound is conservative.

Set $R=0.032$.  For a proposed previous center, \texttt{testBucket} first
tries \eqref{seed:box-support}, and then, if applicable, the dual bound.
It returns success only when the resulting upper bound in
\eqref{seed:distance-bound} is below the safely lowered comparison
threshold for $(NR)^2$.  Hence success certifies the entire feasible
portion of the current bin pair, rather than just its midpoint.
Distance cutoffs, bucket indices, and reach tests only decide which previous
centers to try.  Omitting a possible previous center cannot produce a
successful coverage test.

For clarity, the shared column is counted once in
\eqref{seed:nine-coordinates}, although its displacement enters both
equations in \eqref{seed:stacked-system}.  Reuse requires equality of the
selected \emph{position} $s$.  Moving that position to the start of the
five $X$ coordinates then applies the same coordinate permutation to the
current point and the retained center, and preserves their Euclidean
distance.  The resulting order is the selected $X$ entry, the other four
$X$ entries in order, and the four $Y$ tail entries.  These are precisely
the nine entries of the dephased first two rows and columns used in the
next stage.  Centers retained earlier for the same $X$ tuple are handled
identically, with the first five components of $k$ zero.

\begin{proposition}[Exhaustive seed cover]
\label{seed:cover}
For a completed execution of the specified \texttt{covout} traversal
under the arithmetic conditions of Section~\ref{sec:rounding}, the written
rational midpoint centers, each with its recorded selected position, cover
the seed coordinates of every exact complex Hadamard matrix up to the
equivalences of Lemma~\ref{seed:equivalences}, by closed Euclidean balls
of radius $0.032$.  After the indicated column ordering and conversion of
the centers to binary phase values, radius $0.0320001$ is sufficient.
\end{proposition}

\begin{proof}
Given an exact Hadamard matrix, first use
Lemma~\ref{seed:arrangement} to arrange $x$.  Choose a nondecreasing closed
bin assignment, determine $s$ from that tuple, and then arrange and assign
$y$ with the same shared bin.  All these choices exist also on bin
boundaries and at repeated phases.  The tuple is present in the stated
enumeration.  Lemmas~\ref{seed:gap-filters}, \ref{seed:sum-test}, and
\ref{seed:shared-interval}, together with their certified outward numerical
margins, ensure that its bin pair is not rejected by a necessary filter.

Consider the retained centers as the remaining bin pairs are traversed.
If a coverage test succeeds, Lemma~\ref{seed:dual-support} and the certified
comparison margins put every feasible point of the current pair in a
radius-$R$ ball about an already retained center.  If no test succeeds,
the program retains the current midpoint itself.  Its distance from any
point of its full nine-dimensional bin is at most
\[
 \frac{\sqrt9}{2N}=\frac{3}{224}<0.032.
\]
Thus, by induction over the traversal, the retained centers cover every
feasible point in every processed pair.  Restricting the attempted reuse
centers can only cause further midpoint insertions.  Previously retained
centers remain retained and are included in the output.

The common selected-position permutation preserves all these distances.
Finally, Section~\ref{sec:rounding} bounds the joint error of the nine
binary midpoint phases by less than $10^{-14}$.  The triangle inequality
therefore permits the enlarged radius $0.0320001$.  The centers need not
themselves satisfy the exact Hadamard equations; their balls cover the
necessary seed coordinates of all exact solutions.
\end{proof}
